\section{The grouping principle}\label{sec:theory}
This section makes the picture of \cref{fig:branches} precise. Operators become polynomials (\cref{sec:ring}); the
factors of the polynomial split every solution, and terms do not (\cref{sec:factor}); the split is as stable as the angle
between the factors' solution spaces (\cref{sec:cond}); and finite data can support but not certify the result
(\cref{sec:certify}). Each subsection ends with the rule it gives to the algorithm. Proofs are in \SMproofs.

\subsection{Derivatives become polynomials}\label{sec:ring}
Every operator we consider has constant coefficients,
$P(\partial_t,\partial_x)=\sum_{j,k}c_{jk}\,\partial_t^j\partial_x^k$ with $c_{jk}\in\C$. The exponentials $\psi_{\tau,\xi}(t,x)=e^{\tau t+\xi x}$ are common eigenfunctions of all such operators:
\begin{equation}\label{eq:eigen}
  P(\partial)\,\psi_{\tau,\xi}=P(\tau,\xi)\,\psi_{\tau,\xi}.
\end{equation}
Differentiation becomes multiplication, and operators become polynomials.

\begin{definition}[Symbol and characteristic variety]\label{def:ring}
The constant-coefficient operators form the commutative ring $\ring=\C[\partial_t,\partial_x]\cong\C[\tau,\xi]$. The
polynomial $P(\tau,\xi)$ is the \emph{symbol} of $P(\partial)$, and its zero set
$\Var(P)=\{(\tau,\xi)\in\C^2:P(\tau,\xi)=0\}$ is the \emph{characteristic variety}.
\end{definition}

By \eqref{eq:eigen}, $\psi_{\tau,\xi}$ solves $P(\partial)u=0$ exactly when $(\tau,\xi)\in\Var(P)$. A Fourier mode in
space is the case $\xi=i\kappa$ with wavenumber $\kappa\in\R$; then $P(\lambda,i\kappa)=0$ is the \emph{dispersion
relation}, and its roots $\lambda(\kappa)$ are the \emph{temporal rates} of the modes $e^{\lambda(\kappa)t+i\kappa x}$.
The real part of a rate is growth or decay, the imaginary part is oscillation. We keep $\xi$ in statements about the
ring and set $\xi=i\kappa$ only on Fourier modes, so that a real equation has a symbol with real coefficients.

\begin{example}[Running example]\label{ex:running}
$P=(\partial_t+c\,\partial_x)(\partial_t-\nu\,\partial_x^2)$ has symbol $P=Q_1Q_2$ with $Q_1=\tau+c\,\xi$ and
$Q_2=\tau-\nu\,\xi^2$. On Fourier modes the two factors give two branches of rates, $\lambda_1(\kappa)=-ic\kappa$ (a
wave that travels without decay) and $\lambda_2(\kappa)=-\nu\kappa^2$ (a standing mode that decays). A solution is any
superposition of modes on the two branches. The equation is one object; its two factors are two groups of modes.
\end{example}

\subsection{Groups are factors, not terms}\label{sec:factor}
Let the symbol be factored into pairwise coprime factors,
\begin{equation}\label{eq:factored}
  P=Q_1\cdots Q_N,\qquad Q_j=P_j^{e_j},
\end{equation}
with $P_j$ irreducible and $e_j$ its multiplicity. Two ways of splitting a solution suggest themselves: one piece per
factor $Q_j$, or one piece per additive term of the operator (for $u_t=\nu u_{xx}-cu_x$, one piece for diffusion and one
for transport). Only the first works.

\begin{theorem}[Factor decomposition]\label{thm:factor}
Let $\Omega\subset\R^2$ be bounded, connected and simply connected, and let $P$ be as in \eqref{eq:factored}. Every
solution $u$ of $P(\partial)u=0$ on $\Omega$ (a distribution or a smooth function) can be written
\begin{equation}\label{eq:split}
  u=\sum_{j=1}^Nu_j,\qquad Q_j(\partial)\,u_j=0 .
\end{equation}
The tuples $(u_1,\dots,u_N)$ that sum to zero form a space of exponential-polynomials of dimension exactly
\begin{equation}\label{eq:delta}
  \deltaP=\sum_{i<j}\dim_\C\ring/(Q_i,Q_j),
\end{equation}
the number of intersection points of the varieties $\Var(Q_i)$, counted pair by pair and with multiplicity.
\end{theorem}

The theorem says that the decomposition \emph{exists} and that its only freedom has finite dimension. Each meeting point
of two curves $\Var(Q_i)$ and $\Var(Q_j)$ is a mode that solves both factors and can therefore be moved from one group to
the other without changing $u$. In one variable the statement is B\'ezout's identity: from $aQ_1+bQ_2=1$
follows the split $u=bQ_2u+aQ_1u$. In two variables coprime factors no longer generate $1$; they only have finitely many common
zeros, and $\deltaP$ counts exactly that defect. The hypotheses matter: on a domain with a hole the Laplacian factors as
$(\partial_x+i\partial_y)(\partial_x-i\partial_y)$, yet its solution $\log\abs{z}^2$ on an annulus is not a sum of a
holomorphic and an antiholomorphic single-valued function. This is why the algorithm works on rectangles.

\begin{example}[Running example: two groups, two shared modes]\label{ex:delta}
$Q_1=\tau+c\xi$ and $Q_2=\tau-\nu\xi^2$ meet at $(0,0)$ and $(c^2/\nu,-c/\nu)$, so $\deltaP=2$: a constant and the
exponential $e^{(c^2/\nu)t-(c/\nu)x}$ can be moved between the groups. Among Fourier modes only the first is visible:
the $\kappa=0$ mode belongs to both groups.
\end{example}

\paragraph{The groups are read from the rates} When every factor is an evolution factor, $Q_j=\tau-\Lambda_j(\xi)$, the
component $u_j$ is obtained from $u$ by a Fourier multiplier. On a wavenumber $\kappa$ at which no two branches meet,
\begin{equation}\label{eq:lagrange}
  \hat u_j(t,\kappa)=\prod_{k\neq j}\frac{\partial_t-\Lambda_k(i\kappa)}{\Lambda_j(i\kappa)-\Lambda_k(i\kappa)}\;\hat u(t,\kappa),
\end{equation}
the Lagrange formula for the spectral projectors of $\partial_t$ restricted to that wavenumber. \emph{The rates decide
the groups}, and the poles of the multiplier sit exactly at the visible part of the overlap $\deltaP$. This is what the
algorithm exploits: it never needs the operator in advance, only the rates and their arrangement on curves.

\paragraph{Why additive terms do not form groups}
\begin{proposition}[Terms do not give components]\label{prop:terms}
If $P=A_1+A_2$ with constant coefficients, then every solution of $P(\partial)u=0$ that is a sum $a+b$ with $A_1a=0$
and $A_2b=0$ also solves $A_1A_2u=0$. If $\Var(A_1)\cap\Var(A_2)$ is finite, these solutions form a finite-dimensional
space of exponential-polynomials.
\end{proposition}
\begin{proof}[Proof sketch]
$A_1A_2(a+b)=A_2A_1a+A_1A_2b=0$ because the operators commute. A common zero of $P$ and $A_1A_2$ is a common zero of $A_1$
and $A_2$, so the joint kernel of $P$ and $A_1A_2$ lives on a finite set of exponents.
\end{proof}

For the heat equation, $A_1=\partial_t$ and $A_2=-\nu\partial_x^2$: a ``stationary part plus a diffusion part''
$f(x)+a(t)+b(t)x$ is killed by $\partial_t\partial_x^2$, while a single heat mode $e^{-\nu\kappa^2t}\cos\kappa x$ is not. A
generic solution is not a sum of single-term pieces. What a term-wise split leaves over is the \emph{interaction} of the
terms; it is as large as the solution itself, and on a non-periodic window it cannot be removed by any continuous,
hence any regularized, procedure (\SMinteraction). Different factors, in contrast, superpose exactly: \emph{factors, not
terms, are the groups.}

\roleinalg{The objects to cluster are modes, and a cluster is a factor: a curve $\Var(Q_j)$ in the plane of rates. The
overlap $\deltaP$ is the irreducible ambiguity; the algorithm resolves it by a fixed convention (\cref{sec:group1}).
Because the interaction of terms cannot be regularized away, the algorithm uses no coefficient penalty.}

\subsection{How well separated are the groups? The Friedrichs angle}\label{sec:cond}
\Cref{thm:factor} says that the split exists, not that it is stable. Let $M_j=\ker_{\Hs}Q_j$ be the solutions of factor
$j$ in the Sobolev space $\Hs(\Omega)$, let $\sumop(u_1,\dots,u_N)=\sum_ju_j$ be the sum operator, and define its
conditioning constant
\begin{equation}\label{eq:gammaN}
  \gN=\inf_{0\neq v\perp\ker\sumop}\frac{\norm{\sumop v}^2}{\norm{v}^2},\qquad \norm{(u_1,\dots,u_N)}^2=\sum_j\norm{u_j}_{\Hs}^2 .
\end{equation}
It is the worst case, over all tuples of pieces, of the size of the sum relative to the size of the pieces. If the
spaces are orthogonal, $\gN=1$; if $M_2$ approaches $M_1$, one can take $u_2\approx-u_1$, the pieces stay of unit size
while their sum vanishes, and $\gN\to0$. For two groups $\gamma_2=1-\cF$, where $\cF=\cos\thetaF$ is the cosine of the
smallest nonzero angle between $M_1$ and $M_2$, the \emph{Friedrichs angle}~\cite{Deu01} (\cref{fig:friedrichs}). In
regression terms $M_1$ and $M_2$ are two blocks of regressors, and $\cF$ is the largest correlation between a combination
of the first block and one of the second.

\begin{figure}[tbp]
\centering
\begin{tikzpicture}[scale=1.0,>=Stealth]
  \draw[->,gray] (-0.2,0)--(3.8,0); \draw[->,gray] (0,-0.2)--(0,2.6);
  \draw[very thick,->,color=MidnightBlue] (0,0)--(3.4,0.55) node[right]{\small $M_1=\ker Q_1$};
  \draw[very thick,->,color=BrickRed] (0,0)--(2.0,2.3) node[above]{\small $M_2=\ker Q_2$};
  \draw (1.0,0.16) arc (9.2:49:1.0); \node at (1.25,0.62) {\small$\thetaF$};
  \node[anchor=west,align=left,font=\small] at (5.2,1.3)
   {$\cF=\cos\thetaF$, \quad $\gamma_2=1-\cF$\\[2pt]
    error amplification $\gamma_2^{-1/2}$\\[2pt]
    running example: $\thetaF\approx0.60$, $\gamma_2\approx0.17$};
\end{tikzpicture}
\caption{The Friedrichs angle between the solution spaces of two factors. A small angle means that large pieces can
almost cancel, so recovering the pieces from their sum amplifies errors by up to $\gamma_2^{-1/2}=(1-\cos\thetaF)^{-1/2}$.
The values are those measured on the running example (\cref{sec:worked}).}\label{fig:friedrichs}
\end{figure}

\begin{theorem}[Stability of the split]\label{thm:hm}
Every $u\in\ker_{\Hs}P$ has components $u_j\in\Hs$ if and only if $\sum_jM_j$ is closed and dense in $\ker_{\Hs}P$.
Closedness is equivalent to $\gN>0$, and then the minimum-norm components satisfy
\begin{equation}\label{eq:hm}
  \Big(\sum_j\norm{u_j}_{\Hs}^2\Big)^{1/2}\le\gN^{-1/2}\,\norm{u}_{\Hs}.
\end{equation}
For the finite-dimensional sampled kernels used by the algorithm the sum is always closed, and stability reduces to
$\gN>0$.
\end{theorem}

Two remarks keep the theorem honest. On infinite-dimensional component spaces $\gN>0$ can fail, and the pieces of a
smooth $u$ may be one derivative rougher than $u$. For $N\ge3$ groups the pairwise angles do not determine $\gN$: three
spaces can be pairwise well separated and jointly nearly dependent, so the full Gram operator of $\sumop$ is needed.

\roleinalg{$\gN$ is computed on the sampled kernels and reported with every result. It sets the width of the intervals of
the components (\cref{prop:bayes-angle}): closely aligned groups are never hidden; they show up as wide intervals,
not as a confident but wrong split.}

\subsection{What finite data can and cannot certify}\label{sec:certify}
When does a solution determine its equation? Let the \emph{annihilator} $\Ann(u)=\{Q\in\ring:Q(\partial)u=0\}$ collect
every constant-coefficient equation that $u$ satisfies; it is an ideal of $\ring$. The question is whether
$\Ann(u)=(P)$, i.e.\ whether every equation $u$ satisfies is a consequence of $P$. The answer turns on one class of
functions.

\begin{definition}[Exponential-polynomials]\label{def:ep}
$\EP$ is the set of finite sums $\sum_jp_j(t,x)\,e^{\tau_jt+\xi_jx}$ with polynomials $p_j$ and exponents
$(\tau_j,\xi_j)\in\C^2$. Equivalently (\SMproofs), $w\in\EP$ if and only if the derivatives of $w$ span a
finite-dimensional space.
\end{definition}

An exponential-polynomial \emph{remembers too little}: it sits on finitely many points of $\C^2$, and through finitely
many points pass infinitely many curves, so infinitely many equations. A single heat mode $e^{-\nu\kappa^2t+i\kappa x}$
solves the heat equation, but equally the pure decay $\partial_t+\nu\kappa^2$ and every other equation whose variety
passes through $(-\nu\kappa^2,i\kappa)$. The heat kernel $t^{-1/2}e^{-x^2/(4\nu t)}$, whose derivatives span an
infinite-dimensional space, fills the whole curve, and every equation it satisfies is a multiple of $\partial_t-\nu\partial_x^2$.
The general statement strips one irreducible factor at a time.

\begin{theorem}[Identification]\label{thm:ident}
Let $\Omega$ be open and connected, $P(\partial)u=0$, and $P=\prod_iP_i^{e_i}$ with distinct irreducible $P_i$. Then
$\Ann(u)=(P)$ if and only if $(P/P_i)(\partial)u\notin\EP(\Omega)$ for every $i$. In that case $P$ is the unique
annihilator of minimal degree, up to a constant.
\end{theorem}

In words: $u$ determines its equation exactly when, after removing any one factor, what remains is still a genuine
solution of that factor, not a finite sum of modes. Data, however, are finitely many numbers, and finitely many numbers
are always matched by an exponential-polynomial.

\begin{proposition}[Finite data]\label{prop:certify}
Let $\Omega$ be connected and $z_0\in\Omega$. (i) If the derivatives of $u$ span a space of dimension at most $n$, the
derivatives of order at most $2n-1$ at $z_0$ determine $\Ann(u)$. (ii) Without such a bound, for every finite set of
derivatives at a point, and for every finite set of samples, there is an exponential-polynomial that matches them; hence
no finite data prove $\Ann(u)=(P)$.
\end{proposition}

\roleinalg{An algorithm can refute a candidate equation, or report that it is compatible with the data at the available
resolution, but it can never certify it. The output therefore carries a label (\cref{sec:labels}), and the label that
plays the role of ``identified'' requires each factor to be seen at enough distinct wavenumbers to pin it down within
its class (\cref{prop:counting}).}

